\chapter{Introduction}
\label{ch:intro}

\section{Problem Statement}
Multi-agent systems --- where several autonomous agents interact to achieve goals
individually or collectively --- underpin a growing share of modern technology.
Autonomous vehicles share roads, distributed sensors coordinate surveillance, and
energy grids balance supply across many producers. Unlike traditional control
systems with a single centralized decision maker, these settings contain
independent agents that must learn to coexist, coordinate, or compete. The
difficulty is not merely making each agent individually competent; it is
understanding what emerges when many \emph{learning} agents interact.

Reinforcement learning (RL) is a paradigm in which an agent learns to act by
interacting with an environment and receiving reward feedback. Extending RL to
many simultaneously learning agents --- Multi-Agent Reinforcement Learning
(MARL) --- changes the problem fundamentally. Each agent treats the others as part
of its environment, yet those ``parts'' are themselves learning and changing:
today's optimal response can fail tomorrow once the others adapt. This
\emph{non-stationarity} breaks the assumptions that make single-agent RL
convergent and stable.

Crucially, the \emph{relationship} between agents varies. Sometimes cooperation
benefits everyone; at other times genuine conflict arises, and an agent can
exploit a shared resource for private gain while pushing the cost onto the
collective. Knowing \emph{when} cooperation emerges versus \emph{when}
exploitation dominates is central to deploying autonomous agents responsibly.

This dissertation studies cooperation in MARL using \emph{traffic signal control}
as a representative, inherently cooperative case study. Each intersection is an
agent; because network effects mean that gridlock hurts everyone (including the
intersection that caused it), traffic sits at the \emph{pure-coordination} end of
the cooperation spectrum and is an ideal testbed for isolating the value of
explicit coordination mechanisms.

\section{Research Gap}
Despite rapid progress, several questions remain underexplored:

\begin{itemize}[leftmargin=1.4em]
  \item \textbf{The middle of the spectrum.} Most work targets either purely
  cooperative settings or purely competitive games (chess, poker). The continuum
  between them --- where coordination and competition coexist --- receives less
  systematic study.
  \item \textbf{Reward structure.} While reward shaping is known to matter,
  controlled comparisons of \emph{individual} versus \emph{shared} objectives
  across standardized environments are rare.
  \item \textbf{Fairness.} Algorithms typically optimize total system performance
  without examining how benefits distribute across agents.
  \item \textbf{Scalability of coordination.} How the \emph{value} of a
  coordination mechanism changes as a system grows --- from a handful of agents to
  large networks --- has not been systematically characterized.
\end{itemize}

This work contributes to the first and fourth gaps directly. It asks a deceptively
simple question on a controlled traffic benchmark: \emph{does an explicitly
cooperative algorithm (MAPPO, with a centralized critic) actually outperform a set
of independent learners (IPPO) when the two are matched in every other respect?}
The answer --- and the reason behind it --- forms the spine of this dissertation.
